%%%%%%%%%%%%%%%%%%%%%%%%%%%%%%%%%%%%%%%%%%%%%%%%%%%%%%%%%
%% © DRL CSAIL MIT 2022 - All RIGHTS RESERVED.
%% AUTHORS: RAMIN HASANI and ALEXANDER AMINI
%%
%% License: Creative Commons license
%% Attribution-ShareAlike 4.0 International (CC BY-SA 4.0)
%%%%%%%%%%%%%%%%%%%%%%%%%%%%%%%%%%%%%%%%%%%%%%%%%%%%%%%%%

\documentclass{MITcsail}
\usepackage[numbers]{natbib}
\usepackage{amsthm}
\usepackage{graphicx}
\usepackage{booktabs}
\usepackage{textcomp}   % \textperthousand
\usepackage{amsmath}

\newcommand{\permil}{\textperthousand}

\title{COS 568 Final Project: Task 2, Milestone 3}

\author{Anderson Lee}

\begin{document}

\maketitle
\thispagestyle{firstpagestyle}

\section{Overview}

The goal of the final milestone was to improve the Dynamic PGM--LIPP hybrid on
the two mixed workloads required by the assignment: 10\% insertion / 90\%
lookup and 90\% insertion / 10\% lookup. Dynamic PGM has much better update
behavior than LIPP, but LIPP gives the fastest direct lookup path on the
bulk-loaded data. Our Milestone 2 hybrid inserted new keys into a Dynamic PGM
buffer and periodically flushed them into LIPP, but every lookup still had to
pay for a Dynamic PGM probe before falling back to LIPP. This was especially
bad for negative lookups and for keys that were already in the bulk-loaded
LIPP index.

Our final design, \textsc{BloomHybridPGMLipp}, keeps the same correctness
structure but adds a Bloom-filter fast path over the Dynamic PGM buffer. LIPP
stores the initial bulk-loaded keys. Foreground inserts go to Dynamic PGM and
are also added to the Bloom filter. During lookup, the Bloom filter decides
whether the key could be in the Dynamic PGM buffer. If the Bloom test is
negative, we skip Dynamic PGM and go directly to LIPP; if it is positive, we
probe Dynamic PGM first and then fall back to LIPP on a miss. This keeps the
cheap insertion path while removing most unnecessary Dynamic PGM misses.

\section{Implementation}

The implementation is single-threaded and synchronous. We intentionally removed
the asynchronous flushing machinery explored earlier because it added mutexes,
snapshot indexes, shared-pointer traffic, and background flush interference to
the common lookup path. In contrast, the final lookup path is:
\[
\text{Bloom check} \rightarrow
\begin{cases}
\text{Dynamic PGM lookup, then LIPP fallback}, & \text{if Bloom positive},\\
\text{LIPP lookup only}, & \text{if Bloom negative}.
\end{cases}
\]
The insert path is deliberately shorter:
\[
\text{insert into Dynamic PGM} \rightarrow
\text{add key to Bloom filter} \rightarrow
\text{increment buffer count}.
\]
If the buffered-key count reaches the configured flush threshold, the same
foreground thread performs a synchronous migration: it scans every key--value
pair currently in Dynamic PGM, inserts each pair into LIPP, clears the Dynamic
PGM buffer, resets the buffered-key count, and rebuilds an empty Bloom filter
for the next batch. Until that flush happens, new keys are served from Dynamic
PGM; after the flush, they are served from LIPP. Thus insertion normally avoids
LIPP's more expensive online insert path while still preserving a simple
single-owner state transition at flush time.

The Bloom filter uses splitmix64 hashing with double hashing. Its bit vector is
rounded to a power of two, so each probe can use a bit mask rather than a
modulo. The core state is only the Dynamic PGM buffer, the Bloom filter, the
LIPP index, and counters for the initial size, buffered-key count, and flush
threshold. When the buffered-key count reaches the threshold, the flush loop
iterates over every Dynamic PGM entry, inserts it into LIPP, resets Dynamic PGM,
and rebuilds an empty Bloom filter.

The benchmark variant fields are:
\[
(\text{search method},\ \epsilon,\ \text{flush threshold } \permil,\ \text{Bloom bits/key},\ \text{Bloom hashes}).
\]
Table~\ref{tab:config-search} lists the Bloom-hybrid configurations searched in
the final benchmark. We swept smaller Bloom filters for lookup-heavy workloads
because every lookup pays the Bloom-check cost. For insert-heavy workloads, we
used a focused two-point sweep with larger flush thresholds so most inserts
could remain in Dynamic PGM during the run.

\section{Experimental Setup}

We evaluated FB, Books, and OSMC on the two required mixed workloads. Each
workload contains 2M operations, uses no range queries, and uses a 50\%
negative lookup ratio among lookup operations. Each reported throughput is the
mean of three runs from the same benchmark session. For indexes with
hyperparameters, we selected the row with the highest mean mixed throughput.
The baselines are Dynamic PGM, LIPP, and the Milestone 2 naive hybrid.

\input{analysis_results/final_config_search_table.tex}
\input{analysis_results/final_hyperparams_table.tex}

\section{Results}

Figure~\ref{fig:fb-results} through Figure~\ref{fig:osmc-results} show the 12
required plots: throughput and index size for both workloads on all three
datasets. The Bloom hybrid is the highest-throughput index in all six
dataset/workload pairs. On the lookup-heavy workload, it slightly exceeds LIPP
on FB and Books and substantially exceeds LIPP on OSMC. On insert-heavy
workloads, it is much faster than both LIPP and Dynamic PGM because most
insertions are absorbed by Dynamic PGM while Bloom gating keeps lookup overhead
small.

\begin{figure*}[t]
\centering
\begin{minipage}{0.48\textwidth}
\centering
\includegraphics[width=\linewidth]{analysis_results/final_fb_10pct_insert_90pct_lookup_throughput.png}
\end{minipage}
\hfill
\begin{minipage}{0.48\textwidth}
\centering
\includegraphics[width=\linewidth]{analysis_results/final_fb_10pct_insert_90pct_lookup_size.png}
\end{minipage}

\vspace{0.5em}
\begin{minipage}{0.48\textwidth}
\centering
\includegraphics[width=\linewidth]{analysis_results/final_fb_90pct_insert_10pct_lookup_throughput.png}
\end{minipage}
\hfill
\begin{minipage}{0.48\textwidth}
\centering
\includegraphics[width=\linewidth]{analysis_results/final_fb_90pct_insert_10pct_lookup_size.png}
\end{minipage}
\caption{FB results. Top: 10\% insert / 90\% lookup. Bottom: 90\% insert / 10\% lookup.}
\label{fig:fb-results}
\end{figure*}

\begin{figure*}[t]
\centering
\begin{minipage}{0.48\textwidth}
\centering
\includegraphics[width=\linewidth]{analysis_results/final_books_10pct_insert_90pct_lookup_throughput.png}
\end{minipage}
\hfill
\begin{minipage}{0.48\textwidth}
\centering
\includegraphics[width=\linewidth]{analysis_results/final_books_10pct_insert_90pct_lookup_size.png}
\end{minipage}

\vspace{0.5em}
\begin{minipage}{0.48\textwidth}
\centering
\includegraphics[width=\linewidth]{analysis_results/final_books_90pct_insert_10pct_lookup_throughput.png}
\end{minipage}
\hfill
\begin{minipage}{0.48\textwidth}
\centering
\includegraphics[width=\linewidth]{analysis_results/final_books_90pct_insert_10pct_lookup_size.png}
\end{minipage}
\caption{Books results. Top: 10\% insert / 90\% lookup. Bottom: 90\% insert / 10\% lookup.}
\label{fig:books-results}
\end{figure*}

\begin{figure*}[t]
\centering
\begin{minipage}{0.48\textwidth}
\centering
\includegraphics[width=\linewidth]{analysis_results/final_osmc_10pct_insert_90pct_lookup_throughput.png}
\end{minipage}
\hfill
\begin{minipage}{0.48\textwidth}
\centering
\includegraphics[width=\linewidth]{analysis_results/final_osmc_10pct_insert_90pct_lookup_size.png}
\end{minipage}

\vspace{0.5em}
\begin{minipage}{0.48\textwidth}
\centering
\includegraphics[width=\linewidth]{analysis_results/final_osmc_90pct_insert_10pct_lookup_throughput.png}
\end{minipage}
\hfill
\begin{minipage}{0.48\textwidth}
\centering
\includegraphics[width=\linewidth]{analysis_results/final_osmc_90pct_insert_10pct_lookup_size.png}
\end{minipage}
\caption{OSMC results. Top: 10\% insert / 90\% lookup. Bottom: 90\% insert / 10\% lookup.}
\label{fig:osmc-results}
\end{figure*}

\section{Discussion and Correctness}

The lookup-heavy result depends on making the hybrid lookup path close to LIPP's
lookup path. Most negative lookups and bulk-loaded-key lookups are Bloom
negative, so they avoid Dynamic PGM entirely. Bloom false positives are safe:
they only cause an extra Dynamic PGM miss before checking LIPP. The filter is
small enough that this extra check is cheaper than always probing Dynamic PGM,
which explains why the Bloom hybrid beats the naive hybrid by a large margin.
The lookup-heavy sweep in Table~\ref{tab:config-search} shows that the best
settings are not simply the largest Bloom filters. Configurations around
4 bits/key with 2--3 hashes performed best overall, suggesting that reducing
per-lookup Bloom work was more important than minimizing false positives.

For insert-heavy workloads, the best configuration uses a 20\permil{} flush
threshold with $\epsilon=128$ and a 4-bit, 3-hash Bloom filter. This keeps
foreground insertions in Dynamic PGM for most of the 2M-operation workload and
avoids repeated expensive LIPP insertions. The 10\permil{} insert-heavy setting
was consistently worse, which indicates that flushing too early destroys much of
the insertion advantage. The memory footprint is still close to LIPP because
the bulk-loaded keys remain in LIPP; the extra Dynamic PGM and Bloom structures
are small compared with the LIPP index.

The main correctness invariant is that every inserted key is always either in
Dynamic PGM or in LIPP. Before a flush, each inserted key is placed in Dynamic
PGM and added to the Bloom filter. A Bloom filter has no false negatives for
keys that have been inserted into it, so a lookup for a buffered key will probe
Dynamic PGM. During a flush, every Dynamic PGM entry is inserted into LIPP
before Dynamic PGM and the Bloom filter are reset. Therefore, no keys are
discarded or skipped to improve throughput.

\end{document}
